\begin{abstract}
As artificial intelligence models scale beyond fast-memory capacity, inference latency is increasingly governed by coordination across hierarchical memory systems rather than computation alone. This is especially acute in industrial cyber--physical systems (CPS), where inference must be not merely fast on average but deterministically bounded within PLC scan cycles and TSN time windows. Here we show that such coordination is not continuously optimizable, but instead exhibits intrinsic, regime-dependent limits. We decompose end-to-end latency into four components---compute, locality, inter-tier transfer, and synchronization---and identify three operational regimes---capacity-limited, coordination-dominated, and I/O-limited---defined by two dimensionless ratios: the capacity ratio $\alpha$ and the I/O pressure ratio $\beta$. Transitions between regimes are abrupt rather than gradual.

Using a minimal multi-tier transfer model, we derive analytically interpretable regime boundaries that explain why orchestration strategies effective in coordination-dominated regimes fail predictably once capacity or bandwidth limits are crossed. We build ORION (previously RAMSES), a lightweight orchestrator that coordinates GPU VRAM, host DRAM, and NVMe under regime-aware control. On industrial language and vision workloads, ORION reduces model load time by up to 60\%, inference latency by 35\%, and peak VRAM usage by 15\%, while improving throughput per watt by 18.5\%. Performance remained stable through 72 hours of stress testing, with out-of-SLA yield falling almost tenfold. The timing violation rate of $P_{\mathrm{timing}} = 0.004$ clears IEC~61508 SIL-2, connecting regime-aware memory orchestration to industrial safety compliance. Together, these results establish regime-dependent orchestration as a general principle for memory-bound AI inference and other hierarchical-memory systems, showing that hierarchical memory coordination is regime-dependent in a way that can be written down formally---and once it is, the engineering problem becomes tractable. The analysis reframes optimization as a problem of regime-aware design rather than continuous tuning, and the boundary values we derive give system designers concrete targets for stable, energy-efficient inference within industrial information-system stacks.
\end{abstract}


\begin{IEEEkeywords}
Large Language Models, Memory Optimization, Latency Reduction, GPU Resource Management, System Framework
\end{IEEEkeywords}

